\section{Pembentukan \textit{Spatial Weight Matrix}}
Berdasarkan matriks \textit{similarity} yang telah diperoleh, kemudian disusun \textit{Spatial Weight Matrix} (SWM) seperti pada matriks \ref{swm}  sebagai komponen penting dalam model STARMA. SWM menunjukkan kekuatan pengaruh antar grid, dan tidak hanya mempertimbangkan kedekatan geografis, tetapi juga kemiripan spasio-temporal. Dalam penelitian ini, dibangun tiga variasi bobot sebagai pembanding, yaitu bobot seragam yang memberikan bobot sama untuk seluruh tetangga yang dapat dilihat pada persamaan \ref{swm_uniform}, bobot invers jarak yang menggunakan jarak antar lokasi sebagai dasar pembobotan yang dapat dilihat pada persamaan \ref{swm_inverse}, bobot berdasarkan korelasi silang antar grid yang dapat dilihat pada persamaan \ref{swm_korelasi}. Masing-masing variasi digunakan untuk mengukur efektivitas pembobotan spasial dalam meningkatkan performa prediksi model.